\documentclass[a4paper,11pt]{article}
\usepackage[margin=25mm]{geometry}
\usepackage{fontspec,fvextra}
\setmainfont{STIXGeneral}\setmonofont{Menlo}[Scale=0.78]
\title{Complete computations for the mix criteria}
\author{Human Compression project}\date{Author edition; pending independent review}
\begin{document}\maketitle
\section*{Roles of the complete files}
The GAP file supplies the finite certificate for the literal eleven-cell and
simplex markings used in the coprime-label obstruction. Its group arithmetic is
a retained computational premise. The subsequent product argument is written
in full in the application.

The Python file is bounded validation of the complete six-coordinate argument.
It checks all input flag intersections and actual automorphisms, and witnesses
the failure in their mix. The displayed coordinate proof establishes that
example independently of an extrapolation from test cases. Neither program
proves the general path characterization or all source statements. The replay
driver and the two-declaration Lean source are inventoried separately.

\section*{Required finite certificate}
\VerbatimInput[fontsize=\scriptsize,breaklines=true,breakanywhere=true,breaksymbolleft={},breaksymbolright={},breakanywheresymbolpre={},breakanywheresymbolpost={}]{check_coprime.g}
\section*{Bounded validation of the six-coordinate construction}
\VerbatimInput[fontsize=\scriptsize,breaklines=true,breakanywhere=true,breaksymbolleft={},breaksymbolright={},breakanywheresymbolpre={},breakanywheresymbolpost={}]{check_single_base.py}
\end{document}
